% Danh mục từ viết tắt: giải nghĩa theo chức năng trong đề tài, không dịch chữ theo từng từ.
\chapter*{DANH MỤC TỪ VIẾT TẮT}
\addcontentsline{toc}{chapter}{Danh mục từ viết tắt}
\begin{longtable}{@{}p{.115\linewidth}p{.345\linewidth}p{.485\linewidth}@{}}
\toprule
\textbf{Ký hiệu} & \textbf{Thuật ngữ đầy đủ} & \textbf{Ý nghĩa trong nghiên cứu}\\
\midrule
\endfirsthead
\toprule
\textbf{Ký hiệu} & \textbf{Thuật ngữ đầy đủ} & \textbf{Ý nghĩa trong nghiên cứu}\\
\midrule
\endhead
IDS & Intrusion Detection System & Hệ thống giám sát và phát hiện hoạt động có dấu hiệu xâm nhập.\\
NIDS & Network Intrusion Detection System & Hệ thống IDS phân tích dữ liệu lưu lượng mạng để nhận biết hành vi tấn công.\\
ML & Machine Learning & Phương pháp học quy luật từ dữ liệu nhằm dự đoán nhãn của lưu lượng mới.\\
RF & Random Forest & Bộ phân loại kết hợp nhiều cây quyết định để dự đoán nhãn hoặc xác suất tấn công.\\
LR & Logistic Regression & Bộ phân loại ước lượng xác suất thuộc lớp Attack từ các đặc trưng đầu vào.\\
DT & Decision Tree & Mô hình phân loại dựa trên các nhánh điều kiện đối với giá trị đặc trưng.\\
OOF & Out-of-Fold & Dự đoán trên phần dữ liệu không dùng để huấn luyện mô hình ở lượt chia fold tương ứng.\\
CV & Cross-Validation & Phương pháp chia tập huấn luyện thành nhiều phần để ước lượng chất lượng khi phát triển mô hình.\\
FS & Feature Selection & Quá trình chọn tập đặc trưng đầu vào nhằm cân bằng thông tin dự đoán và chi phí tính toán.\\
EDA & Exploratory Data Analysis & Khảo sát phân bố, chất lượng dữ liệu và quan hệ giữa các đặc trưng trước huấn luyện.\\
TP & True Positive & Số mẫu Attack được mô hình dự đoán đúng là Attack.\\
TN & True Negative & Số mẫu Normal được mô hình dự đoán đúng là Normal.\\
FP & False Positive & Số mẫu Normal bị nhận nhầm là Attack (cảnh báo giả).\\
FN & False Negative & Số mẫu Attack bị nhận nhầm là Normal (tấn công bị bỏ sót).\\
FPR & False Positive Rate & Tỷ lệ Normal bị cảnh báo nhầm, tính bằng FP/(FP+TN).\\
FNR & False Negative Rate & Tỷ lệ Attack bị bỏ sót, tính bằng FN/(FN+TP).\\
F1 & F1-score & Trung bình điều hòa của Precision và Recall; phản ánh chất lượng cân bằng giữa hai đại lượng.\\
MCC & Matthews Correlation Coefficient & Chỉ số tương quan giữa nhãn thật và dự đoán, sử dụng đủ bốn ô ma trận nhầm lẫn.\\
CI & Confidence Interval & Khoảng giá trị ước lượng thể hiện bất định của một đại lượng thống kê theo quy trình đã chọn.\\
AP & Average Precision & Chỉ số tóm tắt quan hệ Precision--Recall theo các ngưỡng dự đoán.\\
SMD & Standardized Mean Difference & Mức chênh lệch trung bình giữa hai nhóm, quy đổi theo độ phân tán của đặc trưng.\\
\bottomrule
\end{longtable}
